%----------------------------------------------------------------------------------------
% Tailored CV — Wissenschaftlicher Mitarbeiter, Dermato-Onkologie
% LMU Klinikum, Prof. Lucie Heinzerling
% Focus: predictive test for immunotherapy response, microbiome-tumour interaction
%----------------------------------------------------------------------------------------

\documentclass[10pt,a4paper,sans]{moderncv}

\usepackage[utf8]{inputenc}
\usepackage[T1]{fontenc}

\moderncvstyle{classic}
\moderncvcolor{blue}

\usepackage[scale=0.78]{geometry}
\usepackage{ragged2e}

%----------------------------------------------------------------------------------------

\firstname{Kundai Farai}
\familyname{Sachikonye}

\address{Auf der Lichtung 19}{82178 Puchheim, Germany}
\mobile{(+49) 1626386344}
\email{kundai.sachikonye@bitspark.com}
\homepage{github.com/fullscreen-triangle}{github.com/fullscreen-triangle}

%----------------------------------------------------------------------------------------

\begin{document}

\makecvtitle

%----------------------------------------------------------------------------------------
%	EDUCATION
%----------------------------------------------------------------------------------------

\section{Education}

\cventry{2019--2021}{PhD candidate, Computational Lipidomics}{Technische Universität München}{}{}{
  Mass spectrometry-based metabolomics, multi-omics statistical modelling, federated
  learning for distributed clinical analysis. Departed to pursue independent research.
}
\cventry{2018--2019}{MSc Computational Biology}{Universität Tübingen}{}{}{}
\cventry{2014--2017}{MSc Pharmaceutical Biotechnology}{HAW Hamburg}{}{}{}
\cventry{2011--2014}{BSc Biochemistry and Cell Biology}{Jacobs University Bremen}{}{}{}

%----------------------------------------------------------------------------------------
%	RESEARCH OUTPUT
%----------------------------------------------------------------------------------------

\section{Research Output}

\cvitem{2025}{
  \textbf{On the Thermodynamic Consequences of Categorical Completion on Biological
  Systems.}
  Structural model of adaptive immune selectivity: MHC binding affinity correlates
  with parent protein categorical richness ($r=0.73$, $p<10^{-12}$, $n=2{,}847$
  IEDB peptides). AUC 0.81 vs NetMHCpan 0.68. Generates patient-level predictions
  about checkpoint inhibitor response: which antigens will be presented, which
  checkpoint is the operative target, and which microbiome-associated proteins
  alter the presentation landscape. \textit{Manuscript.}
}

\cvitem{2025}{
  \textbf{On Disease Epistemology in Blind Receiver.}
  Loop-holonomy self-consistency as unique viable immunotherapy response classifier.
  AUC$=1.00$; 25/25 validation experiments pass. Sparse $\ell_1$ LP for minimum-side-effect
  therapeutic combination selection (checkpoint inhibitor + microbiome modulation
  + adjuvant therapy). \textit{Manuscript. AIMe Registry / TUM Weihenstephan.}
}

\cvitem{2025}{
  \textbf{Syndrome: Categorical Resolution of Disease Trajectories.}
  Disease vector $\mathbf{D}\in[0,1]^8$ integrating genetic expression oscillator
  ($D_G$, tumour transcriptomics), membrane potential oscillator ($D_M$, microbiome
  barrier effects), and calcium oscillator ($D_{Ca}$, immune signalling state).
  Sparse LP for minimum-intervention immunotherapy + microbiome combination selection.
  \textit{Manuscript.}
}

%----------------------------------------------------------------------------------------
%	RESEARCH SOFTWARE
%----------------------------------------------------------------------------------------

\section{Research Software}

\cvitem{lavoisier}{
  \textbf{Metabolomics and Lipidomics Mass Spectrometry Pipeline.}
  Complete pipeline from raw mzML to annotated metabolite table: MS1/MS2 ion
  annotation, spectral library matching (NIST, LIPID MAPS, LipidBlast), feature
  extraction. Dual numerical + CNN visual pipeline (PyTorch); 98.9\,\% NIST accuracy;
  Ray distributed computing; $>$100\,GB mzML support.
  Applicable to microbiome metabolite panels (SCFAs, bile acids, indoles, polyamines)
  and skin tumour lipidomic profiling.
  Virtual instrument:
  \href{https://lavoisier.vercel.app/experiment}{lavoisier.vercel.app/experiment}.
  \href{https://github.com/fullscreen-triangle/lavoisier}{github.com/fullscreen-triangle/lavoisier}
}

\cvitem{borgia}{
  \textbf{First-Principles Molecular Spectroscopy for Metabolite Characterisation.}
  Derives spectroscopic properties of metabolites from bounded phase space geometry
  --- no empirical parameters, no reference standard required. The Triple Equivalence
  Theorem (oscillatory dynamics $=$ categorical state enumeration $=$ partition
  operations) enables characterisation of novel microbial metabolites without
  calibration curves. Validated against NIST for nine elements and six molecules
  with zero free parameters; self-consistency errors $<$1.63\%.
  Live: \href{https://honjo-masamune.vercel.app/tools}{honjo-masamune.vercel.app/tools}.
  \href{https://github.com/fullscreen-triangle/borgia}{github.com/fullscreen-triangle/borgia}
}

\cvitem{gospel}{
  \textbf{Tumour Genomics, Metagenomics, and RNA-seq Integration.}
  $10^6$ variants/second VCF processing; RNA-seq quantification; Bayesian network
  inference, Gaussian mixture models. Applicable to tumour neoantigen profiling,
  somatic variant calling in melanoma, and microbiome metagenomics sequencing.
  Validated: 1000 Genomes Phase~3, gnomAD v3.1.2, UK Biobank.
  Live: \href{https://vivid-symbolism.vercel.app/search}{vivid-symbolism.vercel.app/search}.
  \href{https://github.com/fullscreen-triangle/gospel}{github.com/fullscreen-triangle/gospel}
}

\cvitem{crown-prince}{
  \textbf{Geometric Pharmacology and Tumour Cell Spectral Diagnostics.}
  Derives drug-target binding as $K_d = \exp(-\Delta M \cdot \ln b)$ from partition
  geometry; Hill equation as limiting case. Cell Spectral Hologram: tumour cell input
  $\to$ IR/Raman/fluorescence superposition $\to$ loop-holonomy disease diagnostic
  (AUC\,$=$\,1.00) $\to$ sparse $\eta^*$ therapeutic perturbation vector.
  39/39 NIST binding validations; zero free parameters; 50\,MB vs DrugBank 12\,GB.
  Live: \href{https://crown-prince-four-courners.vercel.app/tools/spectrum}{crown-prince-four-courners.vercel.app/tools/spectrum}.
}

\cvitem{syndrome}{
  \textbf{Immunopeptidology Tools and Combination Therapy Optimisation.}
  Four live tools at \href{https://syndrome-seven.vercel.app}{syndrome-seven.vercel.app}
  operating via 3-channel DFT spectral embedding (36D cosine similarity $\rho$,
  $O(cL\log L)$, no database):
  MHC binding predictor ($\rho_\text{MHC}$, default threshold 0.72);
  neoantigen identifier (dual criterion: $\Delta\rho_\text{self}>0.15$ for immune
  detection AND $\rho_\text{MHC}>0.72$ for presentation, scored as
  $\rho_\text{MHC}\times\Delta\rho_\text{self}$);
  immune escape predictor ($E(v)=-\max\rho(v,\mathrm{Ab})+w\cdot\rho(v,\mathrm{rec})$);
  antibiotic resistance predictor (drug--enzyme spectral complementarity).
  Disease vector $\mathbf{D}\in[0,1]^8$; sparse $\ell_1$ LP for
  minimum-intervention checkpoint inhibitor $+$ microbiome combination selection.
  \href{https://github.com/fullscreen-triangle/syndrome}{github.com/fullscreen-triangle/syndrome}
}

%----------------------------------------------------------------------------------------
%	RESEARCH EXPERIENCE
%----------------------------------------------------------------------------------------

\section{Research Experience}

\cventry{2021--present}{Independent Computational Biology Developer}{Bitspark GmbH}{}{}{
  Multi-omics frameworks, first-principles molecular spectroscopy, immunotherapy
  response modelling. All systems publicly deployed and validated.
}

\cventry{2019--2021}{Computational Lipidomics}{TUM / Bayerisches Forschungsinstitut}{Munich}{}{
  Mass spectrometry pipeline for lipid metabolite characterisation: peak detection,
  ion annotation (LIPID MAPS, LipidBlast), feature engineering, multi-omics
  statistical modelling, federated learning for distributed clinical analysis.
  Python, Java. Instrumentation: TIMS-PASEF, QTOF, ion trap.
}

\cventry{2017--2018}{Glycomics and Proteomics}{Universitätsklinikum Hamburg-Eppendorf}{}{}{
  Automated LC-MS/MS and MALDI-TOF characterisation pipeline for proteins and
  glycans from biological samples; protein identification and quantification.
}

\cventry{2019}{Therapeutic Lipid Production}{Fraunhofer Institut IGB}{Stuttgart}{}{
  Dynamic and linear programming methods for optimisation of therapeutic lipid
  production in industrial photobioreactors. Python, C.
}

\cventry{2013}{Cancer Biomarker Discovery}{University Medical Center Groningen}{}{}{
  UPLC-MS/MS shotgun proteomics for cancer biomarker discovery from serum.
  Protocol optimisation and statistical analysis.
}

%----------------------------------------------------------------------------------------
%	TECHNICAL SKILLS
%----------------------------------------------------------------------------------------

\section{Technical Skills}

\cvitem{Metabolomics / Lipidomics}{
  LC-MS/MS, MALDI-TOF, UPLC-MS/MS, TIMS-PASEF, QTOF; metabolite annotation
  (NIST, HMDB, LIPID MAPS, LipidBlast); SCFA, bile acid, indole, and lipid
  species identification; spectral library matching; feature extraction from
  mzML/mzXML
}

\cvitem{Immunology / Oncology}{
  MHC-I/II peptide presentation modelling, neoantigen prediction, checkpoint
  inhibitor response prediction, tumour microenvironment analysis, pathway analysis
  (KEGG, Reactome, GSEA), single-cell RNA-seq (Scanpy/Seurat)
}

\cvitem{Microbiome Analytics}{
  Metagenomics pipeline (sequence alignment, taxonomic classification, variant
  calling), microbial metabolite panel design, microbiome-host interaction
  multi-omics integration
}

\cvitem{ML / Statistics}{
  PyTorch, scikit-learn, XGBoost; Bayesian networks, Gaussian mixture models,
  variational Bayes; survival analysis; multi-omics integration (MOFA, UMAP)
}

\cvitem{Languages}{Python, R, Rust, TypeScript, Java, C, SQL, Bash}

\cvitem{Infrastructure}{
  Ray, Dask, Docker, FastAPI, AWS, PostgreSQL, memory-mapped I/O for large datasets
}

%----------------------------------------------------------------------------------------
%	LANGUAGES
%----------------------------------------------------------------------------------------

\section{Languages}

\cvitemwithcomment{English}{Native / Fluent}{}
\cvitemwithcomment{German}{Conversational}{Living and working in Germany since 2011}

%----------------------------------------------------------------------------------------

\renewcommand{\listitemsymbol}{-~}

\end{document}
